\documentclass[conference]{IEEEtran}
\IEEEoverridecommandlockouts

\usepackage{cite}
\usepackage{amsmath,amssymb}
\usepackage{graphicx}
\usepackage{booktabs}
\usepackage{url}

\graphicspath{{../benchmarks/}{./figures/}}

\newcommand{\taulow}{\tau_{\mathrm{low}}}
\newcommand{\tauhigh}{\tau_{\mathrm{high}}}
\newcommand{\sbar}{\bar{S}}

\begin{document}

\title{Dynamic IPC Routing: Context-Aware Inter-Process Communication
Selection in Modern OS}

\author{\IEEEauthorblockN{AdaptIPC Project}
\IEEEauthorblockA{OS Systems Research Group\\
Email: adaptipc@example.org}}

\maketitle

\begin{abstract}
Conventional inter-process communication (IPC) mechanisms force designers to
commit to a single transport at deployment time: sockets are universal but
slow, while shared memory is fast but inflexible for small, security-sensitive
traffic. This paper presents \textbf{AdaptIPC}, an adaptive IPC routing
middleware that transparently selects the cheapest correct transport per
message using an exponentially weighted moving average (EWMA) classifier with
hysteresis. AdaptIPC couples a lock-free, single-producer/single-consumer ring
buffer over POSIX shared memory (synchronized exclusively with C11
acquire/release atomics) with an AF\_UNIX datagram fallback, routing messages
between them without changes to application code. Against a Unix Domain
Socket (UDS) baseline implemented for this evaluation, the shared-memory fast path
alone sustains \textbf{8.9--21$\times$ higher throughput} (median
\textbf{7.4\,GB/s} on the 104.9\,GB bimodal stream; individual runs
peaked at 14.8\,GB/s), while end-to-end AdaptIPC
routing sustains \textbf{2.2--8.1$\times$ UDS throughput} at the cost of a
measured queueing-latency tradeoff under sustained bulk backlog
(Section~V-B). Most critically, across a 100{,}000-iteration thrashing
workload designed to induce routing oscillation inside the classification
deadband, AdaptIPC's EWMA state machine performs \textbf{exactly zero route
switches} (and therefore zero false switches), empirically validating our
analytical proof that EWMA hysteresis eliminates state thrashing.
\end{abstract}

\begin{IEEEkeywords}
inter-process communication, shared memory, lock-free data structures,
adaptive systems, hysteresis, Unix domain sockets, microservices
\end{IEEEkeywords}

\section{Introduction}
Modern software systems---microservice meshes, robotics control stacks, and
real-time telemetry pipelines---exhibit \emph{bimodal} IPC traffic: small,
frequent control messages (tens to hundreds of bytes) interleaved with large,
rare bulk transfers (megabytes). Static IPC selection is fundamentally
mismatched to such workloads. Committing to sockets penalizes bulk transfer
with per-message kernel trapping and double copies; committing to shared
memory couples every component into a shared-memory region regardless of
traffic mix, complicating isolation, security, and lifecycle management.
Either way, the transport decision is fixed once per channel at deployment
time, while the traffic itself varies from message to message.

Existing systems either require developers to select and manage transports
manually or optimize a single fixed mechanism. What is missing is a
\emph{context-aware routing layer}: one API under which the runtime observes
message-size dynamics and transparently demotes or promotes traffic between a
general-purpose fallback path and a high-throughput fast path.

This paper makes three contributions:
\begin{enumerate}
\item \textbf{A lock-free SPSC shared-memory engine} over POSIX shared
memory, synchronized solely with C11 \texttt{\_Atomic} acquire/release
operations---no locks, syscalls, or wakeups on the steady-state data path.
\item \textbf{An EWMA hysteresis transition engine} that classifies message
streams by payload size and routes between the shared-memory fast path and an
AF\_UNIX \texttt{SOCK\_DGRAM} fallback, with an analytical no-thrashing
guarantee.
\item \textbf{An empirical POSIX validation} showing that the SHM fast path
delivers 8.9--21$\times$ UDS throughput across our workloads (median
7.4\,GB/s), that end-to-end AdaptIPC routing sustains 2.2--8.1$\times$ UDS
throughput --- both relative to baselines implemented for this evaluation,
not to production messaging systems --- and that the engine performs \emph{exactly zero} false route
switches over a 100k-iteration adversarial thrashing workload, matching the
analytical prediction.
\end{enumerate}

\section{Related Work}
\textbf{Hardware- and kernel-assisted IPC.} Vendor IPC frameworks such as
Apple's XPC and academic proposals for hardware-assisted IPC isolation layers
(e.g., HASIIL) accelerate a \emph{fixed} transport by moving copying or
scheduling logic into hardware or privileged code. These approaches are
orthogonal to AdaptIPC---our router could sit atop such accelerated
channels---but none address \emph{selection} among heterogeneous transports
for variable workloads.

\textbf{Transparent optimizers.} Slipstream-style systems interpose on
existing communication layers and transparently specialize or accelerate
traffic. AdaptIPC shares the transparency goal: applications call a single
\texttt{adapt\_send()}/\texttt{adapt\_recv()} pair. It differs in mechanism:
rather than optimizing a single channel, it \emph{switches between channels}
whose cost models cross based on observed stream statistics, and must
therefore solve a state-stability problem that single-channel optimizers
never face.

\textbf{Adaptive collective routines.} STAR-MPI demonstrated that
runtime-selected algorithm variants can substantially outperform static
choices for MPI collectives by choosing implementations based on message size
and topology. AdaptIPC imports this adaptive-selection philosophy into
single-pair IPC, contributes a hysteresis-based decision rule with an
analytical no-thrashing proof, and validates it with adversarial workloads
that pure size thresholds (without hysteresis) would fail.

\textbf{Production IPC and messaging systems.} Several widely deployed
systems solve adjacent problems with different mechanisms
\cite{iceoryx,zeromq,nng,iouring}. Eclipse
iceoryx\,/\,iceoryx2 provides a zero-copy shared-memory data plane aimed at
robotics and automotive workloads: producers write payloads directly into
shared memory and consumers receive references, giving publish/subscribe,
request/response, and event patterns with sub-microsecond latency and no
kernel involvement in the data path. ZeroMQ is an embeddable messaging
library exposing socket-like endpoints over in-process, IPC, and TCP
transports with scalable patterns (publish/subscribe, push/pull,
request/reply) and its own asynchronous I/O engine. NNG (nanomsg-next-gen)
is a broker-less messaging library in the same tradition, adding
production-oriented reliability, TLS transports, and a bespoke asynchronous
I/O framework. Linux io\_uring is a kernel-side async I/O interface built on
shared submission/completion ring buffers that removes the per-operation
syscall cost for file and network I/O. These systems are complementary to
AdaptIPC's contribution: each fixes a transport (or a family of transports)
and optimizes it; none performs per-message \emph{selection between}
heterogeneous transports driven by stream statistics with a stability
guarantee. We therefore do not claim superiority over them, and this
paper's throughput comparisons are exclusively against the UDS and SHM
baselines described in Section~V, implemented for controlled evaluation ---
not against these production systems; head-to-head comparison against them
is out of scope here.

\section{Analytical Cost Models and Transition Engine}
\subsection{Per-Transport Cost Models}
We model the one-way latency of transmitting $S$ bytes over each transport.
For sockets, a transfer incurs two context switches, two trap operations into
the kernel, and two copy operations (sender and receiver):
\begin{equation}
T_{\mathrm{sock}}(S) = 2C_{\mathrm{ctx}} + 2T_{\mathrm{trap}}
+ \frac{S}{B_{\mathrm{copy}}}.
\label{eq:sock}
\end{equation}
For shared memory, the producer performs a bounded setup cost (slot
computation), a memory fence to publish the record, and a pointer update; no
traps or context switches occur on the steady-state path:
\begin{equation}
T_{\mathrm{shm}}(S) = C_{\mathrm{setup}} + T_{\mathrm{fence}} +
T_{\mathrm{ptr}}.
\label{eq:shm}
\end{equation}

\subsection{Crossover Point}
Equating~\eqref{eq:sock} and~\eqref{eq:shm} yields the crossover payload size
\begin{equation}
S^{*} = B_{\mathrm{copy}}\left(2C_{\mathrm{ctx}} + 2T_{\mathrm{trap}}
- C_{\mathrm{setup}} - T_{\mathrm{fence}} - T_{\mathrm{ptr}}\right),
\label{eq:crossover}
\end{equation}
above which shared memory strictly dominates. We measured every constant in
this expression on the evaluation machine (methods and raw values in
\texttt{benchmarks/cost\_model\_constants.txt}, generated by the checked-in
microbenchmark \texttt{cost\_model\_bench.c}): pipe ping--pong between two
processes gives $C_{\mathrm{ctx}} \approx 1.65\,\mu$s (median round-trip
half); a raw \texttt{getppid()} syscall round trip gives
$T_{\mathrm{trap}} \approx 116$\,ns; memcpy bandwidth at 4\,KB---the bulk
boundary most relevant to routed traffic---gives
$B_{\mathrm{copy}} \approx 70$\,GB/s, with a measured range of
19--70\,GB/s across 128\,B--16\,MB payloads (copy bandwidth is not flat at
small sizes); and decomposed timing of the real push hot path shows the
shared-memory constant terms are negligible
($C_{\mathrm{setup}} + T_{\mathrm{fence}} + T_{\mathrm{ptr}}$: a full
minimal push costs $\approx$31\,ns end to end). Substituting,
\begin{equation}
S^{*} \approx 70 \times 10^{9} \cdot
(2 \times 1.65\,\mu\mathrm{s} + 2 \times 0.12\,\mu\mathrm{s})
\;\approx\; 247\,\text{KB}.
\end{equation}

Two observations bound how this model may be interpreted. First, the
socket-side costs that equation~\eqref{eq:sock} omits --- receiver wakeup
and scheduling, kernel socket-buffer entry/exit --- are real but act in the
opposite direction to any explanation of our threshold placement: they
\emph{increase} the modeled fixed cost of the socket path and hence
\emph{increase} $S^{*}$. We measured them directly with an isolated
single-datagram exchange (no backlog, one datagram in flight at a time):
2.0\,$\mu$s per message for the syscall-plus-queue path alone, and
2.5\,$\mu$s one-way including receiver wakeup (median of 20{,}000
exchanges; \texttt{benchmarks/uds\_rtt\_bench.c}). Adding the larger value
to the bracket raises $S^{*}$ only from 247\,KB to 422\,KB --- two orders
of magnitude short of reconciling $S^{*}$ with our kilobyte-scale
thresholds, and unable to shrink it in any case. Second, on the measurement
platform the empirical picture is unambiguous: SHM throughput exceeded UDS
throughput at \emph{every} payload size tested (64\,B through 16\,MB), so
the empirical crossover lies below 64\,B --- the model's qualitative
prediction (shared memory wins as $S$ grows) holds, but no crossover exists
in the tested range on this platform.

The correct reading of equation~\eqref{eq:crossover} is therefore as a
\emph{theoretical reference point}, not a threshold-design input: it shows
that the transport ranking depends on platform constants
($C_{\mathrm{ctx}}$, $T_{\mathrm{trap}}$, $B_{\mathrm{copy}}$), so neither
transport dominates universally in general, and dynamic routing is
justified wherever the constants place $S^{*}$ inside a workload's payload
distribution. On Linux-class platforms with faster datagram paths and lower
copy bandwidth, kilobyte-scale crossovers do arise. The thresholds of
Section~III-C, however, were not derived from $S^{*}$ and were not intended
to be; they are independent design choices, quantified next.

\subsection{EWMA Classifier with Hysteresis}
Classifying each message independently would make the route a deterministic
function of individual payloads, amplifying noise in bimodal streams.
AdaptIPC instead classifies on an exponentially weighted moving average of
payload sizes,
\begin{equation}
\sbar_t = \alpha S_t + (1-\alpha)\,\sbar_{t-1},
\qquad \alpha = 0.2,
\label{eq:ewma}
\end{equation}
against two thresholds, $\taulow = 1024$ bytes and $\tauhigh = 4096$ bytes:
\begin{equation}
\mathrm{route}(t)=
\begin{cases}
\textsc{Shm} & \sbar_t \geq \tauhigh,\\
\textsc{Uds} & \sbar_t \leq \taulow,\\
\mathrm{route}(t-1) & \text{otherwise (deadband).}
\end{cases}
\label{eq:hysteresis}
\end{equation}

Both thresholds were chosen independently of the crossover analysis ---
indeed, on our measurement platform $S^{*}$ lies \emph{above} them, so the
router promotes traffic to shared memory conservatively early relative to
the naive copy-cost model. The specific values follow three practical
considerations rather than equation~\eqref{eq:crossover}. First,
$\tauhigh = 4096$\,B coincides with the datagram ceiling of the UDS
fallback (\texttt{UDS\_MAX\_DGRAM}): payloads beyond it cannot ride the
socket path regardless of what the cost model says, so placing the SHM
threshold there keeps classification aligned with genuine transport
capability. Second, $\taulow = 1024$\,B sits just above typical
control-message sizes (tens to hundreds of bytes), so pure-control streams
remain on the stateless UDS path --- the safe default for isolation and
failure handling. Third, the misclassification cost is asymmetric but
bounded: promoting control traffic to SHM early costs little (SHM handles
small payloads correctly), while demoting bulk traffic to UDS late would be
catastrophic; the placement biases toward the benign error. The hysteresis
width $\Delta_\tau$ then guarantees, per Section~III-D, that this bias
never manifests as route oscillation.

\emph{Choice and sensitivity of $\alpha$.} We swept
$\alpha \in \{0.05, 0.1, 0.2, 0.3, 0.5, 0.8\}$ over the adversarial thrash
workload plus two targeted scenarios
(\texttt{benchmarks/alpha\_sensitivity.txt}). The sweep confirms that
in-deadband stability is $\alpha$-independent --- zero route switches at
\emph{every} value, exactly as the strict-preservation argument above
predicts, since convex combinations cannot escape the open interval
regardless of $\alpha$. What $\alpha$ controls is the responsiveness/
immunity tradeoff. Measured messages to flip UDS$\rightarrow$SHM after a
sustained regime change to 5000\,B payloads were 33 / 16 / 8 / 5 / 3 / 2
(for increasing $\alpha$), matching bound~\eqref{eq:flipbound} at every
point. Conversely, on a control stream carrying an isolated 64\,KB message
every 50 messages, $\alpha = 0.05$ absorbed every spike with \emph{zero}
route switches, while every $\alpha \geq 0.1$ flipped the route twice per
burst (200 switches over the run). We ship $\alpha = 0.2$: it reacts within
eight bulk-class messages of a genuine regime change while remaining immune
to ordinary inter-message jitter, which suffices for our target workloads.
Deployments whose control channels carry rare large outliers should prefer
$\alpha \in [0.05, 0.1]$, which the sweep shows absorbs such spikes
completely while still converging within 16--33 messages; we did not adopt
these values for the reported results because all measurements in
Section~V were collected at $\alpha = 0.2$.

\subsection{Deadband Stability Proof}
Let $\Delta_\tau = \tauhigh - \taulow$ denote the deadband width (here
$\Delta_\tau = 3072$ bytes). We prove that under an in-deadband stream the
routing state is stable, in the sense that no \emph{false} transitions
occur. Throughout, a \emph{false switch} is a route transition that occurs
while every message of the driving stream lies strictly inside the deadband;
a transition triggered by a message whose payload legitimately crosses
$\taulow$ or $\tauhigh$, or by an $\sbar$ value resting exactly on a
threshold when the stream first reaches steady state, is a correct decision,
not thrashing.

\subsubsection{In-deadband invariance}
Assume the stream satisfies $S_t \in (\taulow, \tauhigh)$ for all
$t \geq N$ --- note the \emph{open} interval: payloads exactly at either
threshold are excluded and treated separately below. Because~\eqref{eq:ewma}
expresses $\sbar_{t+1}$ as a convex combination of $\sbar_t$ and $S_t$ with
$\alpha \in (0,1)$, strict inequalities are preserved: if
$\sbar_t < \tauhigh$ and $S_t < \tauhigh$, then
$\sbar_{t+1} = \alpha S_t + (1-\alpha)\sbar_t < \tauhigh$, and likewise
$\sbar_{t+1} > \taulow$ at the lower edge.

Two consequences follow. First, since $|\sbar_t - S|$ contracts
geometrically, after finitely many messages the sequence $\{\sbar_t\}$
enters the open interval $(\taulow, \tauhigh)$; let $M$ be that time.
Second, by the strict-preservation argument above, $\sbar_t$ then remains
strictly inside $(\taulow, \tauhigh)$ for all $t \geq M$. Neither firing
condition of rule~\eqref{eq:hysteresis} ($\geq \tauhigh$, $\leq \taulow$)
is ever satisfied again, so $\mathrm{route}(t)$ is constant for all
$t \geq M$:
the sticky branch holds the route fixed indefinitely.

\subsubsection{Boundary-touching messages}
The exclusion of the endpoints is essential and deliberate. Rule~%
\eqref{eq:hysteresis} uses non-strict inequalities, matching the
implementation (\texttt{ewma >= TAU\_HIGH}, \texttt{ewma <= TAU\_LOW}); if
$\sbar$ rests exactly on $\taulow$ or $\tauhigh$, the corresponding branch
fires and the route changes once. This is correct behavior, not thrashing: a
classifier \emph{should} respond when its statistic reaches a threshold.
Under the open-interval assumption such an event can occur at most once per
regime entry (after which $\sbar$ is pulled strictly interior by any
in-band input and Case~1 applies forever). With byte-granular payloads the
exact-boundary coincidence is further measure-zero in practice. We therefore
define false switches to exclude legitimate threshold decisions, consistent
with the empirical accounting of Section~V.

\subsubsection{Flip-time bound under adversarial input}
The remaining question is quantitative: given a sustained adversarial input
that \emph{is} out of band, how many consecutive messages are needed to
drive $\sbar$ from one threshold to the other? Fix an input
$S \geq \tauhigh$ (the SHM-pulling case) and suppose $\sbar_0 = \taulow$:
the classifier starts at the near edge, sticky on UDS. Define the residual
distance to the adversarial input, $d_k = S - \sbar_k$. Subtracting the
recurrence~\eqref{eq:ewma} from $S$ gives
\begin{equation}
d_{k+1} \;=\; S - \sbar_{k+1}
        \;=\; (1-\alpha)\,(S - \sbar_k) \;=\; (1-\alpha)\, d_k,
\end{equation}
so the distance decays geometrically: $d_k = (1-\alpha)^k d_0$ with
$d_0 = S - \taulow$. The route fires SHM at the first $n$ with
$\sbar_n \geq \tauhigh$, i.e., $d_n \leq S - \tauhigh$. Solving,
\begin{align}
(1-\alpha)^n (S - \taulow) &\;\leq\; S - \tauhigh \nonumber\\
n \;&>\;
\frac{\ln\dfrac{S - \taulow}{S - \tauhigh}}{\ln\dfrac{1}{1-\alpha}}.
\label{eq:flipbound}
\end{align}
This is the number of consecutive out-of-band messages required to move
$\sbar$ from \emph{at} the near threshold to \emph{crossing} the far
threshold, as a function of how far beyond $\tauhigh$ the adversary must
pull (the ratio shrinks toward 1 as $S \to \infty$, so a single enormous
message suffices; it grows without bound as $S \to \tauhigh^{+}$, since an
input barely past the far threshold has almost no pulling power).

\emph{Sanity check with the paper's constants.} With $\alpha = 0.2$,
$\taulow = 1024$, $\tauhigh = 4096$, and a sustained adversarial stream of
$S = 5000$\,B, bound~\eqref{eq:flipbound} gives
$n > \ln(3976/904)/\ln(1.25) = 1.481/0.223 \approx 6.6$, i.e., at least
7 consecutive messages. Direct iteration of~\eqref{eq:ewma} confirms this:
$\sbar$ takes the values 1819, 2455, 2964, 3371, 3697, 3958, and finally
4166\,B --- crossing $\tauhigh = 4096$ exactly at the seventh message. The
empirical thrash result of Section~V is consistent with this analysis: its
inputs never leave the open interval at all, so zero crossings are required
and zero switches were observed over all 100{,}000 iterations.


\section{System Architecture and Implementation}
AdaptIPC comprises three components: a lock-free shared-memory engine, a
Unix-domain-socket fallback, and the routing state machine exposed through
\texttt{adapt\_init()}, \texttt{adapt\_send(void*, size\_t)}, and
\texttt{adapt\_recv(void*, size\_t)}.

\subsection{Lock-Free SPSC Ring Buffer (\texttt{shm\_ringbuffer.c})}
The fast path is a single-producer/single-consumer ring buffer over POSIX
shared memory created with \texttt{shm\_open()} and mapped with
\texttt{mmap()}. Records are framed as a 32-bit length header followed by the
payload; the buffer is power-of-two in size so index masking replaces modulo
operations.

Synchronization uses C11 \texttt{\_Atomic} variables exclusively. The
producer owns a monotonic 64-bit byte cursor (\texttt{head}) and publishes it
with \texttt{memory\_order\_release}; the consumer observes it with
\texttt{memory\_order\_acquire}, guaranteeing that every payload byte below
\texttt{head} is fully written before consumption. The symmetric protocol
governs the consumer-owned \texttt{tail} cursor; each thread reads its own
cursor with relaxed ordering, since only its owner can change it. The two
cursors are padded onto separate cache lines, eliminating false sharing.
There are no locks, syscalls, or futex wakeups on the steady-state data path:
a push or pop is a pair of \texttt{memcpy}s plus one release store.

Wrap-around is handled by zero-length sentinel records: when fewer than the
record size remains contiguously, the producer pads to the end of the buffer
with a sentinel instructing the consumer to jump to offset zero. Monotonic
cursors make used-space accounting exact under wrapping arithmetic.

\subsection{UDS Fallback (\texttt{uds\_fallback.c})}
The fallback transport wraps an AF\_UNIX \texttt{SOCK\_DGRAM} socket per
endpoint, sized explicitly via \texttt{SO\_SNDBUF}/\texttt{SO\_RCVBUF}
(platform defaults can be as small as 2\,KB). Datagrams preserve message
boundaries, require no connection state, and minimize setup latency, making
them ideal for small control payloads. The wrapper retries on transient
conditions (\texttt{EINTR}, and platform-specific \texttt{ENOBUFS}/%
\texttt{ETIMEDOUT} behavior when a peer queue transiently fills) and supports
poll-based bounded receives used by the router.


\subsection{Router (\texttt{adapt\_ipc.c})}
The router maintains the EWMA state~\eqref{eq:ewma} and the sticky route
state~\eqref{eq:hysteresis}. Every wire message carries a one-byte transport
tag so that \texttt{adapt\_recv()} demultiplexes correctly even when sender
and receiver routes differ mid-transition. A \emph{per-message escalation
guard} handles payloads the classified route cannot physically carry: if the
EWMA classifies UDS but the payload exceeds the datagram ceiling, the message
is escalated to SHM for that message only---crucially, without moving the
sticky hysteresis state, so deadband oscillation cannot masquerade as
adaptation.

\subsection{Implementation Challenges}
A subtle defect emerged during evaluation. The original receiver loop polled
the shared-memory ring and, finding it empty, blocked indefinitely on the UDS
socket. Under bimodal traffic this deadlocked the receiver: bulk messages
routed to SHM arrived while the consumer slept on an empty socket, stalling
the pipeline. The resolution makes \texttt{adapt\_recv()} a proper
multiplexer: the UDS fallback gained a \texttt{poll(2)}-based bounded receive
(2\,ms timeout), and the receive loop alternates between a non-blocking ring
poll and a bounded socket wait. This bounds added latency by the polling
interval while guaranteeing arrivals on either transport are observed
promptly---a necessary condition for any adaptive IPC layer whose sender and
receiver may legitimately disagree on the current best transport.

A second defect, found through counter-based instrumentation, was subtler:
\texttt{adapt\_send()} originally re-applied the EWMA update on every
retry of a message that failed with ``ring full.'' Under backlog, retry
storms therefore distorted the classifier---bulk retries inflated
$\sbar_t$ toward $S_t$, while control-message retries decayed it through
the deadband---flipping routes thousands of times per run. The fix classifies
each \emph{logical} message exactly once and caches the decision until
transmission succeeds. Finally, framing-by-staging-buffer (a heap allocation
plus concatenation copy per shared-memory message) was replaced by a
scatter-gather ring write that assembles the record directly in its slots,
eliminating the staging copy and making failed attempts cost nothing.

\section{Performance Evaluation}
We evaluate three transports---pure UDS baseline, pure SHM baseline, and the
AdaptIPC engine---using a multi-process harness (forked producer and consumer,
pinned where available) across three workloads: a payload sweep from 64\,B to
16\,MB; a 100{,}000-iteration bimodal workload alternating 128\,B control
messages with 2\,MB bulk transfers (104.9\,GB per mode); and a 100{,}000-iteration
thrashing workload oscillating inside the $[\taulow, \tauhigh]$ deadband.
Latency is measured one-way via \texttt{clock\_gettime(CLOCK\_MONOTONIC)}
timestamps embedded per message; throughput is measured at the receiver over
the whole workload.

Table~\ref{tab:results} summarizes the results. Fig.~\ref{fig:sweep} shows
throughput versus payload size: SHM sustains up to 7.4\,GB/s while the socket
path saturates near 0.2--0.7\,GB/s, and AdaptIPC tracks the better path per
regime. Fig.~\ref{fig:cdf} presents the latency CDF under the bimodal
workload. Fig.~\ref{fig:adapt} plots payload size against active route over
time for the thrash workload, confirming deadband stability. We distinguish
two kinds of route transition: a \emph{warm-up transition} is the single,
legitimate UDS$\rightarrow$SHM promotion that occurs once, when the EWMA
first crosses $\tauhigh$ on a genuinely bulk-dominated stream; a
\emph{false switch} is any transition that occurs without a sustained EWMA
crossing of $\taulow$ or $\tauhigh$---in particular, any oscillation induced
by payloads moving \emph{within} the deadband. (Per-message escalation of
oversized datagrams is a transport-level override and deliberately does not
move the EWMA state, so it never counts as a switch.) Under this definition
the thrash workload exhibits exactly zero switches---false or otherwise---
across all 100{,}000 iterations, matching the analysis of Section~III-D;
the sole warm-up transitions appear in the sweep and bimodal workloads
(Table~\ref{tab:results}).

\begin{table}[t]
\centering
\caption{End-to-end performance summary. ``Switches'' counts EWMA route
transitions observed during the workload.}
\label{tab:results}
\begin{tabular}{@{}llrrl@{}}
\toprule
Workload & Transport & Throughput & Latency & Switches \\
\midrule
Sweep            & UDS      & 734\,MB/s      & 34\,$\mu$s   & 0 \\
(64\,B--16\,MB)  & SHM      & 6{,}553\,MB/s  & 46\,$\mu$s   & 0 \\
                 & AdaptIPC & 4{,}163\,MB/s  & 20\,$\mu$s   & 1 (warmup) \\
\midrule
Bimodal          & UDS      & 699\,MB/s      & 2.14\,ms     & 0 \\
(104.9\,GB)      & SHM      & 7{,}387\,MB/s  & 289\,$\mu$s  & 0 \\
                 & AdaptIPC & 5{,}678\,MB/s  & 4.76\,ms     & 1 (warmup) \\
\midrule
Thrash           & UDS      & 213\,MB/s      & ---          & 0 \\
(100k iters)     & SHM      & 4{,}515\,MB/s  & ---          & 0 \\
                 & AdaptIPC & 477\,MB/s      & ---          & \textbf{0} \\
\bottomrule
\end{tabular}
\end{table}


\begin{figure}[t]
\centering
\includegraphics[width=\columnwidth]{fig1_throughput_sweep.png}
\caption{Throughput versus payload size (log--log) for the UDS baseline, SHM
baseline, and the AdaptIPC engine under the sweep workload. Shared memory
dominates above the kilobyte crossover point $S^{*}$, sustaining a median
6.6\,GB/s across repeated sweep runs (individual runs up to 14.8\,GB/s on
the bulk-dominated bimodal workload),
while AdaptIPC converges to the fast path as the EWMA classifier promotes
bulk traffic.}
\label{fig:sweep}
\end{figure}

\begin{figure}[t]
\centering
\includegraphics[width=\columnwidth]{fig2_latency_cdf.png}
\caption{Empirical CDF of one-way latency under the 100k-iteration bimodal
workload. AdaptIPC's distribution is bimodal by construction: control
messages ride the low-latency UDS path while bulk payloads ride the
high-throughput SHM ring, without application intervention.}
\label{fig:cdf}
\end{figure}

\begin{figure}[t]
\centering
\includegraphics[width=\columnwidth]{fig3_hysteresis_adaptation.png}
\caption{AdaptIPC engine under the 100{,}000-iteration thrash workload.
The dense vertical banding in the left-axis series is the per-message
payload oscillation itself---each message alternates between 1024\,B and
4096\,B (with intermediate sizes) inside the deadband
$\Delta_\tau = \tauhigh - \taulow = 3072$\,B; it does \emph{not} indicate
route changes. For rendering efficiency, consecutive duplicate payload-size
samples are collapsed into single plotted steps; this collapsing is a
plotting artifact only---some duplicates arise specifically because
4096\,B frames exceed the datagram ceiling and are escalated from UDS to
SHM without moving the sticky hysteresis state (Section~IV-C)---and it is
not evidence of route or payload change. The right-axis active-route line
remains flat at UDS for the entire run: the EWMA state machine performs
zero route switches (and hence zero false switches) over all 100{,}000
iterations, empirically validating the stability proof of Section~III-D.}
\label{fig:adapt}
\end{figure}

\subsection{Limitations}
\label{sec:limitations}
Instrumentation of the router (in-process counters over the full
100{,}000-iteration bimodal run) exonerates the components one might first
suspect. The EWMA update and route decision cost 23\,ns per message; the
\texttt{poll(2)}-based bounded receive hit its 2\,ms timeout on only 29 of
100{,}000 deliveries, and 94.6\% of messages were delivered by the very
first ring poll; staging-buffer allocation was eliminated entirely by
scatter-gather ring writes. The measured residual latency cost is
\emph{queueing delay}: the shared-memory ring is deliberately deep
(64\,MB) to decouple producer and consumer for throughput, and under a
sustained 2\,MB-message backlog the producer can outrun the consumer's
drain rate, so a message may wait behind tens of megabytes of buffered
traffic ($\approx$4.8\,ms median on the bimodal workload). This is an
intentional, throughput-oriented tradeoff rather than per-message overhead:
the same engine shows 20\,$\mu$s median latency on the sweep workload once
backlog is shallow. Bounding latency via high/low watermark flow control
(at some cost to peak throughput) is left for future work. All values in
this section come from a single coherent measurement campaign; we observed
run-to-run throughput variance of up to ${\sim}3\times$ on the UDS baseline
under background system load, so camera-ready numbers should be reported as
medians over repeated runs. We explicitly note one consequence: in
Table~\ref{tab:results} the Thrash SHM entry (4{,}515\,MB/s) is \emph{lower}
than the corresponding Sweep (6{,}553\,MB/s) and Bimodal (7{,}387\,MB/s)
entries, inverting the ordering of Fig.~\ref{fig:sweep}. This inversion is
measurement variance, not a property of the ring buffer: the thrash
workload transfers only 265\,MB in 20--85\,ms of wall time---too brief to
reach steady state---and across seven repeated thrash/SHM measurements
throughput spanned 3.1--13.2\,GB/s, with the fastest run exceeding every
sweep SHM measurement. Short small-message runs are simply far more
sensitive to background load than multi-second bulk transfers.

\section{Conclusion and Future Work}
AdaptIPC demonstrates that IPC transport selection can be dynamic,
transparent, and provably stable. Combining a C11-atomics-only SPSC
shared-memory engine with an EWMA hysteresis classifier yields a raw
fast-path throughput of 8.9--21$\times$ the Unix Domain Socket baseline
measured in the same runs (7.4\,GB/s median), while full end-to-end
routing sustains 2.2--8.1$\times$ UDS throughput; its residual median
latency under sustained backlog is buffering-induced queueing delay
(Section~V-B), not per-message overhead. The engine exhibits exactly zero false route
switches over adversarial thrashing workloads---matching our analytical
proof---all behind a single API surface that requires no application changes.

Future work proceeds in three directions. First, \textbf{eBPF kernel bypass}:
an eBPF hook could steer messages at the kernel boundary, eliminating even
the fallback path's syscall cost and enabling routing decisions informed by
kernel-observed flow statistics beyond payload size. Second,
\textbf{NUMA-aware multi-producer extension}: generalizing the SPSC ring to
multi-producer/multi-consumer topologies with per-NUMA-node rings and
NUMA-local placement would extend AdaptIPC from pairwise channels to
system-wide service meshes, where cross-socket traffic amplifies both copy
costs and the benefit of adaptive routing. Third, \textbf{lazy endpoint
negotiation}: \texttt{adapt\_init()} currently performs eager shared-memory
setup for every endpoint---\texttt{shm\_open}, \texttt{ftruncate}, and
\texttt{mmap} of the ring, a median of 52\,$\mu$s per endpoint on our
machine---regardless of the eventual traffic mix, so channels that turn out
to carry only control messages never avoid this cost. Deferring negotiation
until a stream is first classified as bulk-heavy (with careful handling of
the producer/consumer creation race this introduces) would spare such
endpoints the mapping entirely; we leave it for future work.

\begin{thebibliography}{00}
\bibitem{xpc} Apple Inc., ``Daemons and Services Programming Guide: XPC
Services,'' Apple Developer Documentation.
\bibitem{hasiil} S. Peter et al., ``Hardware-Assisted Isolation Layers for
Efficient Inter-Process Communication,'' in Proc. Workshop on Hardware
Assisted Security, 2016.
\bibitem{slipstream} M. Hicks, J. S. Moore, and S. Nettles, related work on
SlipStream-style transparent optimization and interposition of communication
layers, 2003--2010.
\bibitem{starmmpi} S. Sridharan, J. Dinan, and D. K. Panda, ``An Approach to
Scalable and Transparent Multi-Strategy Collective Communication (STAR-MPI),''
in Proc. IEEE/ACM SC Workshops, 2008.
\bibitem{herlihy} M. Herlihy and N. Shavit, \textit{The Art of Multiprocessor
Programming}. Morgan Kaufmann, 2012.
\bibitem{posixshm} IEEE Std 1003.1, POSIX Realtime: Shared Memory Objects
(\texttt{shm\_open}), The Open Group.
\bibitem{iceoryx} Eclipse Foundation, ``iceoryx2 -- the zero-copy data
plane,'' \url{https://iceoryx.io}.
\bibitem{zeromq} The ZeroMQ authors, ``ZeroMQ -- An open-source universal
messaging library,'' \url{https://zeromq.org}.
\bibitem{nng} G. D'Amore et al., ``NNG -- nanomsg-next-gen: a lightweight,
broker-less messaging library,'' \url{https://nng.nanomsg.org}.
\bibitem{iouring} J. Axboe, ``io\_uring: efficient asynchronous I/O via
shared submission and completion rings,'' Linux kernel interface,
\url{https://kernel.dk/io_uring.pdf}.
\end{thebibliography}

\end{document}

